\documentclass[11pt]{article}
\usepackage[margin=1in]{geometry}
\usepackage{amsmath,amssymb}
\usepackage{booktabs}
\usepackage{enumitem}
\usepackage{hyperref}
\usepackage{xcolor}
\usepackage{listings}

\lstset{
  basicstyle=\ttfamily\footnotesize,
  breaklines=true,
  frame=single,
  numbers=left,
  numberstyle=\tiny\color{gray},
  showstringspaces=false,
  columns=fullflexible,
  keepspaces=true
}

\newcommand{\perc}{\mathrm{perc}}
\newcommand{\self}{\mathrm{self}}
\newcommand{\file}[1]{\texttt{#1}}

\title{Three-way \texttt{percRecip} Diagnosis Log\\[4pt]
\large May 19, 2026}
\author{Jinwoo Cho}
\date{}

\begin{document}
\maketitle

\section{Goal Statistic}

Recover the parameter $\beta$ for
\[
s \;=\; \sum_i \sum_j Y[\self]_{i,j} \cdot Y[i]_{j,i}
\]
across $K$ perceiver slices, with a single shared $\beta$.  Intuition:
``actor $i$'s self-reported tie to $j$ is reciprocated in $i$'s own
perceived view.''

%% ============================================================
\section{Why the Original Intra-Slice Implementation Failed}
%% ============================================================

The original \file{ThreeWayPerceivedReciprocityEffect} treated the
statistic as intra-slice:
\[
s_i \;=\; \sum_k Y[i]_{i,k} \cdot Y[i]_{k,i}
\]

This is \emph{identically zero} in the three-way fork.
\file{initializeFRAN.r:2289--2296} marks perceiver $i$'s own row of
slice $Y[i]$ as structural-zero (value~10):
\begin{lstlisting}[language=R]
## Mark perceiver k's own row as structural zeros (value 10).
## Perceiver k's self-reported ties (x_{k,k,*}) are handled
## exclusively by Y[self]; setting row k to structural 10 prevents
## the simulation engine from creating ministeps in these positions
## and avoids double-counting with Y[self].
for (tt in 1:TT) {
  arr3[k, , tt] <- ifelse(is.na(arr3[k, , tt]), NA, 10)
}
\end{lstlisting}

Hence $Y[i]_{i,k} = 0$ structurally and $s_i \equiv 0$.  C++ debug
(\texttt{[percRecip tieStat total]}) confirmed $> 22$ million calls
with \emph{zero} non-zero returns.  This file is now orphaned;
see \file{src/model/effects/ThreeWayPerceivedReciprocityEffect.h}
header notice.

%% ============================================================
\section{Replacement: Cross-Network Direction-1 Effect}
%% ============================================================

The target statistic is fundamentally cross-network.  In the unpacked
threeway tensor:
\[
Y[\self]_{i,j} \;=\; \text{depvar}[i,i,j]
\quad\text{(self-net diagonals)},
\qquad
Y[i]_{j,i} \;=\; \text{depvar}[i,j,i]
\quad\text{(perceived row $j \neq i$)}.
\]
The natural implementation is on each perceived slice $Y[i]$'s
objective, referencing $Y[\self]$ via \texttt{interaction1}.  For a tie
flip $(e\!g\!o, alter)$ in $Y[i]$ with perceiver index $i$:
\[
\Delta s_i \;=\;
\begin{cases}
Y[\self]_{i, ego} & \text{if } alter = i\\
0                  & \text{otherwise}
\end{cases}
\]

This is implemented as
\file{src/model/effects/generic/PerceiverRestrictedInTieFunction.{h,cpp}},
wired in \file{EffectFactory.cpp} as
\begin{lstlisting}[language=C++]
else if (effectName == "percRecip")
{
    pEffect = new GenericNetworkEffect(pEffectInfo,
        new PerceiverRestrictedInTieFunction(
            pEffectInfo->interactionName1(),
            pEffectInfo->variableName()));
}
\end{lstlisting}
The function parses the perceiver index from the \emph{outcome}
variable name (\texttt{Y[k]} $\to k\!-\!1$) and gates on
\texttt{alter == perceiverIndex} before returning
\texttt{Y[self]\_\{alter, ego\}} via \texttt{InTieFunction}'s
\texttt{pNetworkCache()->inTieValue(alter)}.

%% ============================================================
\section{Open Issue (Blocking Recovery)}
%% ============================================================

Despite the conceptually correct cross-network implementation, MoM
recovery still fails: \texttt{percRecip-hat} stuck at $0$,
\texttt{tconv.max = NA}, all SEs \texttt{NA}.  The same failure
\textbf{occurs with the standard, well-tested \texttt{crprodRecip}
effect under our harness}.  Diagnostic sequence:

\begin{enumerate}[leftmargin=*]
  \item \textbf{\file{diagnose\_percRecip\_simulator.R}} (May 19):
        sim with $\beta\!=\!0$ vs $\beta\!=\!5$ on Y[self] objective
        produced identical wave-2 (stat $=3$ in both, diff $=0$).
  \item Confirmed \textbf{Y[self] does not evolve}: 0/400 cells changed
        across all draws, even with rate=30.  Y[i] does evolve
        (45/8000 cells).  Therefore effects on Y[self]'s objective
        cannot contribute to any choice probability.
  \item Switched to Direction~1 (effect on $Y[i]$, interaction1 =
        $Y[\self]$).  Still $\beta\!=\!0$ vs $\beta\!=\!5$ yields
        identical trajectories.
  \item \textbf{\file{diagnose\_baseline\_crprodRecip.R}}: the
        \emph{standard} \texttt{crprodRecip} also fails to influence
        the simulator under our harness.  $\beta\!=\!0$ vs
        $\beta\!=\!5$ both give stat $=76$.  So the bug is
        \textbf{not} in our custom function.
  \item \textbf{\file{diagnose\_no\_share\_crprodRecip.R}}: turning
        off \texttt{shareParameters} and using a single slice does
        not fix it either.  $\beta$ still has no effect.
  \item \textbf{\file{diagnose\_estimation\_mode.R}}: even in proper
        estimation mode (not \texttt{simOnly}), with R-side manually
        generated wave-2 that has 1583/8000 cells changed and a
        strong embedded \texttt{crprodRecip} signal,
        \texttt{crprodRecip-hat} stays at $0$ with \texttt{tconv.max
        = 0.234} (false ``convergence'').
  \item Adding \texttt{sharedCov=TRUE} to \texttt{sienaDependent}
        (the missing argument from the user's earlier successful
        \file{tests/test\_crprod\_effects.R}) did not change the
        outcome.
\end{enumerate}

\paragraph{Conclusion.}
The cross-network mechanism (any \texttt{crprod}-family effect on a
perceived slice $Y[i]$ with \texttt{interaction1 = "Y[self]"}) does not
currently influence the simulator's choice probabilities for threeway
data in this fork, regardless of \texttt{shareParameters},
\texttt{sharedCov}, share-canonical setup, or
estimation-vs-\texttt{simOnly} mode.  The bug is somewhere in the path
between effect registration (effects table $\to$ requestedEffects
$\to$ C++ factory) and the simulator's evaluation function for $Y[i]$
when the referenced \texttt{interaction1} network is $Y[\self]$.

\paragraph{Next steps to investigate (future session).}
\begin{itemize}[leftmargin=*]
  \item Verify whether \texttt{percX} / \texttt{percSenderAgree}
        (covariate-based three-way effects) recover correctly in this
        harness.  If yes, the issue is specifically with
        cross-NETWORK effects on $Y[i]$; if no, the whole $Y[i]$
        objective machinery is impaired.
  \item Trace the simulator-side wiring: confirm that for slice
        $Y[i]$, the \texttt{percRecip} / \texttt{crprodRecip} effect
        actually ends up in \texttt{NetworkVariable}'s evaluation
        function with the correct $Y[\self]$ cache.
  \item Compare against the user's older
        \file{simulation\_crprod\_recovery.R} on a real-data run to
        confirm whether the supposed earlier ``working'' state really
        recovered cross-network parameters or just converged without
        recovering them.
\end{itemize}

%% ============================================================
\section{Files Touched This Session}
%% ============================================================

\begin{tabular}{@{}lll@{}}
\toprule
File & Status & Purpose \\
\midrule
\file{src/.../PerceiverRestrictedInTieFunction.h} & new   & header for cross-network function\\
\file{src/.../PerceiverRestrictedInTieFunction.cpp} & new & impl with perceiver gating\\
\file{src/.../ThreeWayPerceivedReciprocityEffect.h} & gutted & deprecation stub\\
\file{src/.../ThreeWayPerceivedReciprocityEffect.cpp} & gutted & deprecation stub\\
\file{src/.../AllEffects.h} & edit & drop old percRecip include, add notice\\
\file{src/.../EffectFactory.cpp} & edit & rewire \texttt{percRecip} branch\\
\file{src/sources.list} & edit & swap old percRecip src for new fn src\\
\file{src/Makevars.win} & edit & same as sources.list\\
\file{data/allEffects.csv} & edit & \texttt{percRecip} row $\to$ cross-network\\
\file{tests/recovery\_percRecip.R} & rewrite & moved from project root + new flow\\
\file{tests/diagnose\_*.R} & new (6 files) & step-wise isolation diagnostics\\
\file{tests/recovery\_percRecip\_*.{rds,csv}} & artifacts & last-run outputs\\
\file{notes/percRecip\_diagnosis\_0519.tex} & new & this document\\
\bottomrule
\end{tabular}

\end{document}
